\RequirePackage[utf8]{inputenc}
%\RequirePackage{hgbpdfa}

\documentclass[language=english]{hgbarticle}

\usepackage{authblk}	% for easier formatting of author block
\renewcommand\Authfont{\sffamily\normalsize}
\renewcommand{\Affilfont}{\rmfamily\small}

%%%----------------------------------------------------------
\begin{document}
%%%----------------------------------------------------------

\title{The \textsf{hagenberg-thesis} Package}
\date{\hgbDate}

%% TODO: make sure authors end up in PDF metadata without affiliation!
\author{W.\ Burger}
\author{W.\ Hochleitner}
\affil{University of Applied Sciences Upper Austria\\
		   Department of Digital Media, Hagenberg}

%%%----------------------------------------------------------
\maketitle
%%%----------------------------------------------------------

\begin{abstract}\noindent
The \textsf{hagenberg-thesis} package is a collection of LaTeX
templates for academic theses (bachelor, master, diploma or PhD programs) and
related documents. This manual describes the main features of this package
and how they may be customized.
Pre-configured document templates for English and German manuscripts and a
complete tutorial are available on the package's home repository.
\end{abstract}

\clearpage
\tableofcontents
\clearpage

\section{Introduction}

\makeatletter
\AB@authors
\makeatother

The complete source of this package and auxiliary materials are available on
CTAN%
\footnote{\url{https://ctan.org/pkg/hagenberg-thesis}}
and its development repository.%
\footnote{\url{https://github.com/Digital-Media/HagenbergThesis}}
The package is made available under the terms of the Creative Commons
Attribution 4.0 International Public License.%
\footnote{\url{https://creativecommons.org/licenses/by/4.0/legalcode}}
The only exception are the texts of the detailed documentation of AI tool usage in
\texttt{hgbaiusage-cognisance.sty} (see Sec.~\ref{sec:AIUsage}), which are licensed under the
Creative Commons Attribution-ShareAlike 4.0 International Public License.%
\footnote{\url{https://creativecommons.org/licenses/by-sa/4.0/legalcode}}


\section{Document classes}

The \texttt{hgb} package provides the following document classes, which are
based on the standard \latex\ classes \texttt{book}, \texttt{report}, and
\texttt{article}, respectively:
%
\begin{itemize}
    \item \textbf{\texttt{hgbthesis}} (\texttt{book}):
    for bachelor's, master's, diploma, and PhD theses as well as internship reports,
    \item \textbf{\texttt{hgbreport}} (\texttt{report}):
    for project and term reports,
    \item \textbf{\texttt{hgbarticle}} (\texttt{article}):
    for drafting journal articles.
\end{itemize}


\subsection{Class options}

\subsubsection{General options}

All three document classes accept the following general options:
%
\begin{itemize}
	\item \texttt{language=}\{\texttt{english}|\texttt{german}\}:
		main document language (default: \texttt{english}).%
		\footnote{The bare options \texttt{english} and \texttt{german} still work but are deprecated.
		Note that \texttt{german} implicitly loads Babel's \texttt{ngerman} setup; when switching languages
		(see Sec.~\ref{sec:LanguageSwitchingCommands}) only \texttt{german} should be used.}
	\item \texttt{twoside=}\{\texttt{false}|\texttt{true}\}:
		page layout (default: \texttt{false} = single-sided).
		The bare option \texttt{oneside} is deprecated (use \texttt{twoside=false}).
	\item \texttt{smartquotes=true}:
		smart quotes replacement (\verb!"!...\verb!"! $\rightarrow$ ``...'');
		default is \texttt{smartquotes=false},
	\item \texttt{apa=true}:
		use APA bibliography style instead of \texttt{numeric-comp};
		default is \texttt{apa=false},
	\item \texttt{codestyle=}\{\texttt{color}|\texttt{mono}|\texttt{classic}\}:
		appearance of code listings, see Sec.~\ref{sec:CodeStyle} (default: \texttt{color}),
	\item \texttt{updatecheck=true} (default): perform check of package release date;
		use \texttt{updatecheck=false} to suppress.
\end{itemize}


\subsubsection{Class-specific options}

In addition, the following class-specific options are accepted:

\noindent\textbf{\texttt{hgbthesis}}:
	\begin{itemize}
		\item \texttt{type=}%
			\{\texttt{bachelor}|\texttt{master}|\texttt{diploma}|\texttt{phd}|\texttt{internship}\}:
				thesis type (default: \texttt{master}),
		\item \texttt{titlelanguage=}\{\texttt{english}|\texttt{german}\}:
			language on title pages, defaults to main document language if omitted,
		\item \texttt{theme=}\{\textit{themeID}\}:
			cover page theme style (default: \texttt{default}, see Sec.~\ref{sec:Themes}),
		\item \texttt{degree=}\{\textit{key}\}:
			academic degree on the title page, \eg, \texttt{engineering} (see Sec.~\ref{sec:AcademicDegrees}),
		\item \texttt{proposal=true}: flag to indicate proposal documents (\texttt{proposal=false} is the default).
	\end{itemize}

\noindent\textbf{\texttt{hgbreport}}:
	\begin{itemize}
		\item \texttt{titlepage=}\{\texttt{true}|\texttt{false}\}:
			title page on/off (default: \texttt{true}).
			The bare option \texttt{notitlepage} is deprecated (use \texttt{titlepage=false}).
	\end{itemize}

\noindent\textbf{\texttt{hgbarticle}}:
	\begin{itemize}
		\item \texttt{twocolumn=true}: two-column page layout (default: \texttt{twocolumn=false}),
		\item \texttt{review=true}: review mode with line numbers (default: \texttt{review=false}).
	\end{itemize}

\noindent
For example, to start a master's thesis in German, simply place
%
\begin{LaTeXCode}[numbers=none]
\documentclass[type=master,degree=engineering,language=german,smartquotes=true]{hgbthesis}
\end{LaTeXCode}
%
at the beginning of the document.

The \texttt{proposal} option is intended for a \emph{thesis proposal} (``Exposé'' or ``Proposal'') and is
only effective in \emph{conjunction} with the main thesis types, \eg,
%
\begin{LaTeXCode}[numbers=none]
\documentclass[type=bachelor,language=english,titlelanguage=german,proposal=true]{hgbthesis}
\end{LaTeXCode}
%
The result is a short exposé that contains only one chapter. Thus,
chapter numbers are not displayed. Remove the \texttt{proposal} option to
migrate a proposal document straight to the final thesis (and restore the usual
numbering scheme).

The example also shows that the main document language (\texttt{language=english}) and the
language on the title page (\texttt{titlelanguage=german}) can be set independently.

\subsection{Thesis settings (class \texttt{hgbthesis})}

\texttt{hgbthesis} supports several types of thesis documents. The following
parameters must be specified at the beginning of the main document:
%
\begin{itemize}
    \item \verb!\title{...}!,
    \item \verb!\subtitle{...}! (optional),
    \item \verb!\author{...}!,
    \item \verb!\programtype{...}! (optional, derived from \texttt{type} by default, see Sec.~\ref{sec:AcademicDegrees}), 
    \item \verb!\programname{...}!,
    \item \verb!\institution{...}! (optional if provided by the theme, see Sec.~\ref{sec:ThemeBrand}),
    \item \verb!\placeofstudy{...}!,
    \item \verb!\academicdegree{...}! (optional, free text replacing the degree selected by the class option \texttt{degree}, see Sec.~\ref{sec:AcademicDegrees}),
    \item \verb!\dateofsubmission{yyyy}{mm}{dd}!,
    \item \verb!\advisor[]{...}! (optional, multiple \verb!\advisor! statements allowed),
\end{itemize}
%
Note that \texttt{hgbthesis} only supports a \emph{single author} inside the
\verb!\author{...}! macro argument (commands \verb!\and! and
\verb!\thanks{...}! are deactivated)!

The command \verb!\advisor[role]{name}! accepts an optional argument to associate a \emph{role} 
with the advisor's \emph{name}, \eg, \verb!\advisor[Supervisor]{Prof.~Marie Curie, PhD}!.
Multiple advisors can be specified, for example
%
\begin{LaTeXCode}[numbers=none]
  \advisor[1.~Betreuerin]{Professor Frida K.~Putnik, PhD}
  \advisor[2.~Betreuer]{Franz Grillparzer, TU Wien}
  \advisor[Gutachter]{Dr.~B.\,Rutal, MIT}
\end{LaTeXCode}

\noindent
Other (optional) settings for \texttt{hgbthesis} include:
\begin{itemize}
    \item \verb!\license{cc|strict}!:
    Use the \emph{Creative Commons License} (\texttt{cc} = default) or strict license terms
      (``all rights reserved'').
    \item \verb!\logofile{xxx.pdf}!:
    %Specifies a custom logo image for the title page, to be placed in \verb!images/xxx.pdf!
    %(default is \texttt{logo.pdf}). 
		This command is \textbf{deprecated} -- the thesis logo 
		should be specified in the associated \emph{theme} file instead (see Sec.~\ref{sec:Themes}).
\end{itemize}

\subsection{Degree program type and academic degree}
\label{sec:AcademicDegrees}

The type of the degree program and the academic degree on the title page are derived from the
class options \texttt{type} and \texttt{degree}, so that all theses of a degree program use the
same wording.
The \emph{type of degree program} is derived from \texttt{type} and printed in the title page
language (see Table~\ref{tab:ProgramTypes}). \verb!\programtype{...}! replaces it by any
text, \eg, for an internship report in a master's program.

\begin{table}[H]
\caption{Type of degree program derived from the class option \texttt{type}.}
\label{tab:ProgramTypes}
\centering\small
\begin{tabular}{@{}lll@{}}
	\toprule
	\texttt{type=} & German & English \\
	\midrule
	\texttt{bachelor}   & Bachelorstudiengang & Bachelor's degree program \\
	\texttt{master}     & Masterstudiengang   & Master's degree program \\
	\texttt{diploma}    & Diplomstudiengang   & Diploma degree program \\
	\texttt{phd}        & Doktoratsstudium    & Doctoral program \\
	\texttt{internship} & Bachelorstudiengang & Bachelor's degree program \\
	\bottomrule
\end{tabular}
\end{table}

The \emph{academic degree} is selected by the class option \texttt{degree=\textit{key}}, where
\textit{key} denotes the subject group and the level (bachelor or master) is taken from
\texttt{type}. For example,
%
\begin{LaTeXCode}[numbers=none]
\documentclass[type=master,degree=engineering,...]{hgbthesis}
\end{LaTeXCode}
%
results in ``Master of Science in Engineering (MSc)'' on the title page.
Table~\ref{tab:AcademicDegrees} lists all supported keys. They cover all academic degrees of
Austrian FH degree programs as defined by AQ Austria (board decision of July 7, 2021).%
\footnote{\url{https://www.aq.ac.at/de/akkreditierung/fachhochschulen/downloads.php}}
An unknown key, or a key that is not available for the given type (\eg, \texttt{degree=security-management}
with \texttt{type=bachelor}), results in an error.
The academic degree is shown by themes based on \texttt{default} (see Sec.~\ref{sec:Themes}),
but not for internship reports and not for proposals (\texttt{proposal=true}). Thus the option can
remain unchanged when a proposal is turned into the final thesis.

Academic degrees not covered by the option (\eg, of universities, of other countries, or doctoral
degrees) are specified as free text by \verb!\academicdegree{...}!, for example,
%
\begin{LaTeXCode}[numbers=none]
\academicdegree{Master of Science (MSc)}
\end{LaTeXCode}
%
If both are given, \verb!\academicdegree! takes precedence and a warning is issued.
Custom themes may define additional keys by
%
\begin{LaTeXCode}[numbers=none]
\hgbDefineDegree{<key>}{<bachelor degree>}{<master degree>}
\end{LaTeXCode}
%
where an empty argument means that the degree does not exist at that level.

\begin{table}[H]
\caption{Supported values of the class option \texttt{degree} (academic degrees of Austrian FH degree programs).}
\label{tab:AcademicDegrees}
\centering\small
\begin{tabular}{@{}lll@{}}
	\toprule
	\texttt{degree=} & \texttt{type=} & Academic degree \\
	\midrule
	\texttt{arts-design}         & \texttt{bachelor} & Bachelor of Arts in Arts and Design (BA) \\
	                             & \texttt{master}   & Master of Arts in Arts and Design (MA) \\ \addlinespace
	\texttt{engineering}         & \texttt{bachelor} & Bachelor of Science in Engineering (BSc) \\
	                             & \texttt{master}   & Master of Science in Engineering (MSc) \\ \addlinespace
	\texttt{social-sciences}     & \texttt{bachelor} & Bachelor of Arts in Social Sciences (BA) \\
	                             & \texttt{master}   & Master of Arts in Social Sciences (MA) \\ \addlinespace
	\texttt{business}            & \texttt{bachelor} & Bachelor of Arts in Business (BA) \\
	                             & \texttt{master}   & Master of Arts in Business (MA) \\ \addlinespace
	\texttt{military-leadership} & \texttt{bachelor} & Bachelor of Arts in Military Leadership (BA) \\
	                             & \texttt{master}   & Master of Arts in Military Leadership (MA) \\ \addlinespace
	\texttt{police-leadership}   & \texttt{bachelor} & Bachelor of Arts in Police Leadership (BA) \\ \addlinespace
	\texttt{security-management} & \texttt{master}   & Master of Arts in Security Management (MA) \\ \addlinespace
	\texttt{cultural-studies}    & \texttt{bachelor} & Bachelor of Arts in Cultural Studies (BA) \\
	                             & \texttt{master}   & Master of Arts in Cultural Studies (MA) \\ \addlinespace
	\texttt{natural-sciences}    & \texttt{bachelor} & Bachelor of Science in Natural Sciences (BSc) \\
	                             & \texttt{master}   & Master of Science in Natural Sciences (MSc) \\ \addlinespace
	\texttt{health-studies}      & \texttt{bachelor} & Bachelor of Science in Health Studies (BSc) \\
	                             & \texttt{master}   & Master of Science in Health Studies (MSc) \\ \addlinespace
	\texttt{law}                 & \texttt{bachelor} & Bachelor of Laws (LLB) \\
	                             & \texttt{master}   & Master of Laws (LLM) \\
	\bottomrule
\end{tabular}
\end{table}


\section{Style files and user commands}

The package comes with a set of style (\texttt{*.sty}) files that can be used
independently of the document classes listed above:
\begin{itemize}
  \item \texttt{hgb.sty}: language and date setup, custom commands;
  \item \texttt{hgbabbrev.sty}: various abbreviation macros;
  \item \texttt{hgbaiusage.sty}: documentation of AI tool usage (for \texttt{hgbthesis}; texts of the
    detailed form in \texttt{hgbaiusage-cognisance.sty}, CC BY-SA 4.0);
  \item \texttt{hgbalgo.sty}: additions to \texttt{algpseudocodex} package;
  \item \texttt{hgbbib.sty}: bibliography setup;
  \item \texttt{hgbdict.sty}: language dictionary functions;
  \item \texttt{hgbheadings.sty}: definition of page headers;
  \item \texttt{hgblistings.sty}: setup for code listings;
  \item \texttt{hgbmath.sty}: setup and commands for math typesetting;
  \item \texttt{hgbpdfa.sty}: setup for PDF/A generation.
\end{itemize}

\subsection{General user commands and environments (\texttt{hgb.sty})}
\label{sec:GeneralUserCommands}

\begin{itemize}
   \item \textbf{\texttt{{\bs}hgbDate}}: Outputs the package version date,
    \eg, ``\texttt{\hgbDate}''.
	\item \textbf{\texttt{{\bs}getcurrentlabel}}:
			Yields the most recently assigned label number.
   \item \textbf{\texttt{{\bs}calibrationbox}}\verb!{width}{height}!:
			Inserts a test box for checking the final print size (in millimeters).
	\item \verb!\begin{block}...\end{block}!:
			Dummy environment, provides a limited scope for variable/command redefinitions.
\end{itemize}


\subsection{Language switching commands and smart quotes (\texttt{hgb.sty})}
\label{sec:LanguageSwitchingCommands}

\begin{itemize}
	\item \verb!\SetLanguage{<language>}!:
    Switches to the specified language.
		This and the other language switching commands listed below
		also refresh the \texttt{smartquotes} settings (if activated).
	\item \verb!\begin{SwitchLanguage}{<language>}...\end{SwitchLanguage}!:
    Temporarily switches to the specified language and reverts to the current language 
		afterwards. 
	\item \verb!\begin{english}...\end{english}!:
    Temporarily switches to English language settings, short for
		\verb!\begin{SwitchLanguage}{english}..!.
	\item \verb!\begin{german}...\end{german}!:
    Temporarily switches to German language settings, short for
		\verb!\begin{SwitchLanguage}{german}..!.
	\item \verb!\SmartquotesActive!:
		For debugging only. Returns \texttt{true}/\texttt{false} if smart quotes are activated in
		the current context (or an appropriate message if the \texttt{csquotes} package was not loaded at all).
\end{itemize}


\subsection{Text commands (\texttt{hgbabbrev.sty})}

\subsubsection*{Special characters:}

\begin{itemize}
    \item \textbf{\texttt{{\bs}bs}}: Inserts a backslash character (short for
    \verb!\textbackslash!).
    \item \textbf{\texttt{{\bs}obnh}}: Inserts an optional break with no
    hyphen (\eg, \verb!PlugIn{\obnh}Filter!).
\end{itemize}


\subsubsection*{German abbreviations:}

\begin{itemize}
    \item \textbf{\texttt{{\bs}bzgl}}: 	\bzgl
    \item \textbf{\texttt{{\bs}bzw}}: 	\bzw
    \item \textbf{\texttt{{\bs}ca}}: 		\ca
    \item \textbf{\texttt{{\bs}dah}}: 	\dah
    \item \textbf{\texttt{{\bs}Dah}}: 	\Dah
    \item \textbf{\texttt{{\bs}ds}}: 		\ds
    \item \textbf{\texttt{{\bs}etc}}: 	\etc
    \item \textbf{\texttt{{\bs}evtl}}: 	\evtl
    \item \textbf{\texttt{{\bs}ia}}: 		\ia
		\item \textbf{\texttt{{\bs}oa}}: 		\oa
    \item \textbf{\texttt{{\bs}sa}}: 		\sa
    \item \textbf{\texttt{{\bs}so}}: 		\so
    \item \textbf{\texttt{{\bs}su}}: 		\su
    \item \textbf{\texttt{{\bs}ua}}: 		\ua
    \item \textbf{\texttt{{\bs}Ua}}: 		\Ua
    \item \textbf{\texttt{{\bs}uae}}: 	\uae
    \item \textbf{\texttt{{\bs}usw}}: 	\usw
    \item \textbf{\texttt{{\bs}uva}}: 	\uva
    \item \textbf{\texttt{{\bs}uvm}}: 	\uvm
    \item \textbf{\texttt{{\bs}va}}: 		\va
    \item \textbf{\texttt{{\bs}vgl}}: 	\vgl
    \item \textbf{\texttt{{\bs}zB}}: 		\zB
    \item \textbf{\texttt{{\bs}ZB}}: 		\ZB
\end{itemize}

\subsubsection*{English abbreviations:}

\begin{itemize}
    \item \textbf{\texttt{{\bs}ie}}: 	\ie
    \item \textbf{\texttt{{\bs}eg}}: 	\eg
    \item \textbf{\texttt{{\bs}etc}}: \etc
    \item \textbf{\texttt{{\bs}Eg}}: 	\Eg
    \item \textbf{\texttt{{\bs}wrt}}: \wrt
\end{itemize}

\noindent
Note that none of the above abbreviation macros ``eats'' subsequent white
space, \ie, they can be used without additional controls, as in
``\verb!\wrt what I said!'', for example.

\subsection{Bibliography commands (\texttt{hgbbib.sty})}

\begin{itemize}
    \item 
    \textbf{\texttt{{\bs}AddBibFile}}\newline
    A wrapper to \texttt{biblatex}'s \verb!\addbibresource! macro (for backward compatibility only).
    \item
    \textbf{\texttt{{\bs}MakeBibliography}}[\emph{options}]
    \newline
    Inserts the
    reference section or chapter. By default, references are automatically
    split into category subsections.%
    \footnote{Predefined reference categories are \texttt{literature},
        \texttt{avmedia}, \texttt{online} and \texttt{software}.}
    Use the option \texttt{nosplit} to produce a traditional (\ie,
    contiguous) list of references.
    \item
    \textbf{\texttt{{\bs}mcite}}%
        [\emph{text1}]\{\emph{key1}\}%
        [\emph{text2}]\{\emph{key2}\}%
        \ldots
        [\emph{textN}]\{\emph{keyN}\}%
    \newline
    Analogous to \texttt{bib\-la\-tex}'s \texttt{{\bs}cites} command%
    \footnote{%
    \url{http://mirrors.ctan.org/macros/latex/contrib/biblatex/doc/biblatex.pdf}
    (see Sec.~3.8.3)} but inserts semicolons between reference entries for
    better readability.
\end{itemize}

% \MakeBibliography ... creates a reference section split subsections (default)
% \MakeBibliography[nosplit] ... creates a one-piece reference section


\subsection{Code environments (\texttt{hgblistings.sty})}
\label{sec:CodeEnvironments}

The code environments listed in Table~\ref{tab:CodeEnvironments} are defined. Their keywords correspond
to current language versions (\eg, C23, C++23, C\#~13, Java~21, Python~3, PHP~8.4, Swift~6); built-in
types and functions form a second keyword class. The table also lists the language names for
including code files with \verb!\lstinputlisting!, for example,
%
\begin{LaTeXCode}[numbers=none]
\lstinputlisting[language=Java]{code/Main.java}
\end{LaTeXCode}
%
Languages that are also defined by \texttt{listings} are extended by the dialect \texttt{[hgb]}, which is
set as default dialect where possible (\eg, \texttt{language=Java} selects \texttt{[hgb]Java}). For HTML,
\latex, Scala, SQL and XML, the dialect must be specified explicitly (\eg, \texttt{language=\{[hgb]SQL\}});
otherwise, the original definition of \texttt{listings} is used.

\begin{table}[H]
\caption{Code environments and language names for \texttt{\textbackslash lstinputlisting}.}
\label{tab:CodeEnvironments}
\centering\small
\begin{tabular}{@{}lll@{}}
	\toprule
	Language     & Environment          & \texttt{language=} \\
	\midrule
	Bash         & \texttt{BashCode}    & \texttt{bash} \\
	C            & \texttt{CCode}       & \texttt{C} \\
	C++          & \texttt{CppCode}     & \texttt{C++} \\
	C\#          & \texttt{CsCode}      & \texttt{CSharp} \\
	CSS          & \texttt{CssCode}     & \texttt{CSS} \\
	Dart         & \texttt{DartCode}    & \texttt{Dart} \\
	GLSL         & \texttt{GlslCode}    & \texttt{GLSL} \\
	Go           & \texttt{GoCode}      & \texttt{Go} \\
	HTML         & \texttt{HtmlCode}    & \texttt{\{[hgb]HTML\}} \\
	Java         & \texttt{JavaCode}    & \texttt{Java} \\
	JavaScript   & \texttt{JsCode}      & \texttt{JavaScript} \\
	JSON         & \texttt{JsonCode}    & \texttt{JSON} \\
	Kotlin       & \texttt{KotlinCode}  & \texttt{Kotlin} \\
	\latex       & \texttt{LaTeXCode}   & \texttt{\{[hgb]TeX\}} \\
	MATLAB       & \texttt{MatlabCode}  & \texttt{Matlab} \\
	Objective-C  & \texttt{ObjCCode}    & \texttt{ObjectiveC} \\
	PHP          & \texttt{PhpCode}     & \texttt{PHP} \\
	Python       & \texttt{PythonCode}  & \texttt{Python} \\
	R            & \texttt{RCode}       & \texttt{R} \\
	Rust         & \texttt{RustCode}    & \texttt{Rust} \\
	Scala        & \texttt{ScalaCode}   & \texttt{\{[hgb]Scala\}} \\
	SQL          & \texttt{SqlCode}     & \texttt{\{[hgb]SQL\}} \\
	Swift        & \texttt{SwiftCode}   & \texttt{Swift} \\
	TypeScript   & \texttt{TsCode}      & \texttt{TypeScript} \\
	XML          & \texttt{XmlCode}     & \texttt{\{[hgb]XML\}} \\
	YAML         & \texttt{YamlCode}    & \texttt{YAML} \\
	generic code & \texttt{GenericCode} & \texttt{\{\}} \\
	\bottomrule
\end{tabular}
\end{table}
%
\texttt{hgblistings} is based on the \texttt{listingsutf8}%
\footnote{\url{https://ctan.org/pkg/listingsutf8}}
package, thus any valid \texttt{listings}%
\footnote{\url{https://ctan.org/pkg/listings}}
option may be used; for example, the option \texttt{numbers=none} to suppress
line numbers:
%
\begin{LaTeXCode}[numbers=none]
    \begin{JavaCode}[numbers=none]
    ... // Java code comes here
    \end{JavaCode}
\end{LaTeXCode}


\subsubsection{Code style}
\label{sec:CodeStyle}

The appearance of all code environments is selected by the class option \texttt{codestyle}:
%
\begin{itemize}
	\item \texttt{codestyle=color} (default): white background with a thin gray line between the line
		numbers and the code; syntax highlighting with colors similar to common development environments
		(Visual Studio Code) with sufficient contrast (WCAG~2 AA) for keywords (dark blue), built-in types
		and functions (teal), strings (dark red) and comments (green, italic),
	\item \texttt{codestyle=mono}: the same layout without colors (comments in italics), \eg, for printing
		in black and white,
	\item \texttt{codestyle=classic}: gray background without colors, as in previous versions of this
		package.
\end{itemize}
%
The styles are also defined as \texttt{listings} styles \texttt{hgbcolor}, \texttt{hgbmono} and
\texttt{hgbclassic}, so they can be selected for individual listings as well:
%
\begin{LaTeXCode}[numbers=none]
    \begin{JavaCode}[style=hgbmono]
    ... // Java code without colors
    \end{JavaCode}
\end{LaTeXCode}
%
The following colors are used and may be redefined in the preamble:
%
\begin{itemize}
	\item \texttt{ListingsKeywordColor}: keywords,
	\item \texttt{ListingsTypeColor}: built-in types and functions,
	\item \texttt{ListingsStringColor}: strings,
	\item \texttt{ListingsCommentColor}: comments,
	\item \texttt{ListingsLineNumberColor}: line numbers,
	\item \texttt{ListingsRuleColor}: line between the line numbers and the code,
	\item \texttt{ListingsBackgroundColor}: background (\texttt{classic} only).
\end{itemize}
%
For example,
%
\begin{LaTeXCode}[numbers=none]
    \definecolor{ListingsKeywordColor}{HTML}{0000FF}
\end{LaTeXCode}


\subsection{Mathematical commands (\texttt{hgbmath.sty})}

\texttt{hgbmath} requires (and automatically loads) the \texttt{amsmath}%
\footnote{\url{https://ctan.org/pkg/amsmath}}
package, thus, all commands and symbols of \texttt{amsmath} are available by
default. The following \emph{additional} commands can only be used in math mode:
%
\begin{itemize}
    \item \textbf{\texttt{{\bs}Cpx}}: $\Cpx$ (complex numbers),
    \item \textbf{\texttt{{\bs}N}}: $\N$ (natural numbers),
    \item \textbf{\texttt{{\bs}R}}: $\R$ (real numbers),
    \item \textbf{\texttt{{\bs}Q}}: $\Q$ (rational numbers),
    \item \textbf{\texttt{{\bs}Z}}: $\Z$ (integer numbers).
\end{itemize}


\subsection{Algorithms (\texttt{hgbalgo.sty})}

\texttt{hgbalgo} is a stand-alone package that is based on -- and extends --
the \texttt{algorithmicx} and \texttt{algpseudocodex} packages.%
\footnote{\url{https://ctan.org/pkg/algorithmicx},
    \url{https://ctan.org/pkg/algpseudocodex}}
It fixes some (mostly indentation-related) problems, adds color, and provides
some additional commands. It also loads the \texttt{algorithm}%
\footnote{\url{https://ctan.org/pkg/algorithms}}
package, which defines a compatible float container for algorithms:
\verb!\begin{algorithm}! \verb!...! \verb!\end{algorithm}!.

\paragraph{Additional algorithm commands:}
\begin{itemize}
    \item
    \textbf{\texttt{{\bs}StateNN[<nesting>]\{<text>\}}}:
    Creates a \emph{non-numbered} statement like \texttt{algo\-rith\-micx}'s
    \verb!\Statex! command but provides controlled indentation inside nested
    constructs. The optional integer argument \verb!<nesting>! can be used to
    specify the \emph{nesting depth} to compensate for a bug in
    \texttt{algorithmicx} (the nesting level inside a block is not set
    properly before the first \verb!\State! command). Omitting the optional
    argument should give correct indentation in most situations.
    \item
    \textbf{\texttt{{\bs}Input\{<text>\}}}:
    For describing the input parameters in a procedure's preamble.
    \item
    \textbf{\texttt{{\bs}Output\{<text>\}}}:
    For describing the output values in a procedure's preamble.
    \item
    \textbf{\texttt{{\bs}Returns\{<text>\}}}:
    For describing the return values in a procedure's preamble.
\end{itemize}

\paragraph{Vertical spacing commands:}
The following commands are provided for fine-tuning the vertical spacing
between individual statements of an algorithm (the standard spacing commands
like \verb!\smallskip! \etc\ have no effect between statements):%
\footnote{Note that the standard spacing commands work \emph{between}
\texttt{procedure} and \texttt{function} blocks in the usual way.}
\begin{itemize}
    \item \textbf{{\bs}\texttt{algsmallskip}}: inserts 3pt extra space,
    \item \textbf{{\bs}\texttt{algmedskip}}: inserts 6pt extra space,
    \item \textbf{{\bs}\texttt{algbigskip}}: inserts 12pt extra space.
\end{itemize}
They are supposed to be used inside (\ie, at the end of) statements, for
example:
%
\begin{LaTeXCode}[numbers=none]
    \State $x \gets x + 1$ \algsmallskip
\end{LaTeXCode}

\paragraph{Defined algorithm colors:}
\begin{itemize}
    \item \textbf{\texttt{AlgKeywordColor}}: for algorithm keywords,
    \item \textbf{\texttt{AlgProcedureColor}}: for procedure and function names.
\end{itemize}
These colors can be redefined at any time (see the \texttt{xcolor}%
\footnote{\url{https://ctan.org/pkg/xcolor}}
package), \eg, by
\begin{LaTeXCode}[numbers=none]
    \definecolor{AlgKeywordColor}{named}{black}
    \definecolor{AlgProcedureColor}{rgb}{0.0, 0.5, 0.0}     % dark green
\end{LaTeXCode}


\subsection{Documentation of AI tool usage (\texttt{hgbaiusage.sty})}
\label{sec:AIUsage}

Package \texttt{hgbaiusage} is loaded by \texttt{hgbthesis} and creates the documentation of the
usage of generative AI tools, which follows the affidavit, \ie, it is placed directly after
\verb!\maketitle!. It starts a new unnumbered chapter, and all texts are printed in the title page
language (option \texttt{titlelanguage}). The texts are those of the FH O\"O templates (2026), as is the
text of the affidavit created by \verb!\maketitle!.

\subsubsection*{Short form}

The environment \texttt{aiusage} composes the documentation of predefined text modules, for example,
%
\begin{LaTeXCode}[numbers=none]
\begin{aiusage}
  \aiusedfor{literature-search}
  \aiusedfor{proofreading}
\end{aiusage}
\end{LaTeXCode}
%
The following commands are available inside the environment:
\begin{itemize}
	\item \verb!\aiusednone!: no generative AI tools were used (cannot be combined with the other commands),
	\item \verb!\aiusedfor{<key>}!: AI tools were used for a predefined purpose (see Table~\ref{tab:AIUsagePurposes}),
	\item \verb!\aiusedforother{<text>}!: AI tools were used for another purpose; \verb!<text>! continues
		the sentence like the predefined text modules (see the examples below).
\end{itemize}
An empty environment results in a warning. Examples of other purposes in English and German:
%
\begin{LaTeXCode}[numbers=none]
\aiusedforother{create the figures in Chapter 3.}
\aiusedforother{um die Abbildungen in Kapitel 3 zu erzeugen.}
\end{LaTeXCode}

\begin{table}[H]
\caption{Predefined purposes for \texttt{\bs aiusedfor} (short form).}
\label{tab:AIUsagePurposes}
\centering\small
\begin{tabular}{@{}lp{0.6\linewidth}@{}}
	\toprule
	\textit{key} & Text (English) \\
	\midrule
	\texttt{ideas}                    & gather initial ideas for potential research questions and relevant theories. \\
	\texttt{literature-search}        & search for literature. \\
	\texttt{literature-understanding} & better understand academic literature (e.g., summaries). \\
	\texttt{proofreading}             & proofread parts of the thesis and improve the linguistic style. \\
	\bottomrule
\end{tabular}
\end{table}

\subsubsection*{Detailed form}

The environment \texttt{aiusagedetails} creates a form with two statements on the responsible use
of AI tools (with check boxes) and a table of eleven activities, listing the AI tools used for each
activity and a description of their use. Its texts were adapted by FH O\"O (2026) from
Gimpel et al.\ (2023)%
\footnote{Gimpel, H., Dilger, P., L\"ammermann, L., Urbach, N. (2023). COGNISANCE: Declaration of
Generative AI Tool Usage in Higher Education Tests. Research Center for Information Management.
\url{https://doi.org/10.17605/OSF.IO/VGXF7}}
and, unlike the rest of the package, are licensed under CC BY-SA 4.0. Therefore, they are kept in the
separate file \texttt{hgbaiusage-cognisance.sty}, and the source reference is printed below the form.
The detailed form replaces the short form, since only one of them is intended; if both are used, a
warning is issued. For example,
%
\begin{LaTeXCode}[numbers=none]
\begin{aiusagedetails}
  \aiconfirm{accuracy}
  \aiconfirm{responsibility}
  \aiactivity{coding}{GitHub Copilot}{Generation of test functions (Chapter 4)}
  \aiactivity{editing}{DeepL Write}{Proofreading of Chapters 1--5}
  \aifurtheruse{Generation of test data for the evaluation (Section 5.2).}
\end{aiusagedetails}
\end{LaTeXCode}
%
The environment only collects the following entries; the form is created at its end.
\begin{itemize}
	\item \verb!\aiconfirm{accuracy}!, \verb!\aiconfirm{responsibility}!: checks the box of the
		statement that the outputs of AI tools were verified and corrected where required, and that
		the responsibility lies with the author(s), respectively. A statement that is not confirmed
		results in a warning.
	\item \verb!\aiactivity{<key>}{<tools>}{<description>}!: AI tools used for an activity (see
		Table~\ref{tab:AIUsageActivities}) and a description of their use (type of usage, affected
		sections, \etc). Multiple entries for the same activity are appended; activities without
		entries remain empty (or are omitted with the option \texttt{compact}, see below).
	\item \verb!\aifurtheruse{<text>}!: further use of AI tools (optional; without it, \emph{None} or \emph{Keine} is printed
		below the heading).
\end{itemize}
%
By default, the table lists all activities, as the FH O\"O template does. With the optional argument
\texttt{compact}, only the activities with entries are listed; if no activity has an entry, a
sentence stating this replaces the table:
%
\begin{LaTeXCode}[numbers=none]
\begin{aiusagedetails}[compact]
  ...
\end{aiusagedetails}
\end{LaTeXCode}
%
The sample documents contain an example line for every activity, most of them commented out.

\begin{table}[H]
\caption{Activities for \texttt{\bs aiactivity} (detailed form).}
\label{tab:AIUsageActivities}
\centering\small
\begin{tabular}{@{}ll@{}}
	\toprule
	\textit{key} & Activity (English) \\
	\midrule
	\texttt{ideas}          & Idea generation \& conceptualization \\
	\texttt{literature}     & Literature search and analysis \\
	\texttt{methodology}    & Methodology \\
	\texttt{coding}         & Coding \\
	\texttt{data}           & Data collection and analysis \\
	\texttt{interpretation} & Interpretation and validation \\
	\texttt{structure}      & Structuring and planning the text \\
	\texttt{generation}     & Generating the text \\
	\texttt{translation}    & Translating text \\
	\texttt{editing}        & Reviewing \& editing the text \\
	\texttt{references}     & Citation management \\
	\bottomrule
\end{tabular}
\end{table}


\subsection{PDF/A generation (\texttt{hgbpdfa.sty})}

Package \texttt{hgbpdfa} contains definitions for generating PDF/A-compliant 
(PDF/A-2b) output files. It is based on the emerging \latex-3 \texttt{pdfmanagement} package.%
\footnote{\url{https://ctan.org/pkg/pdfmanagement-testphase}}
\texttt{hgbpdfa.sty} is loaded by \texttt{hgbarticle.cls}. For \texttt{hgbthesis} and
\texttt{hgbreport}, it is loaded at the very beginning of the main document, \ie, before
\verb!\documentclass! (as in all sample documents), and may be omitted if no PDF/A output is
required. The font encoding (\texttt{T1}) and the Latin Modern fonts (\texttt{lmodern}) are loaded
by all classes independently of \texttt{hgbpdfa} when compiling with \texttt{pdflatex}.
PDF metadata insertion is done with the associated \texttt{hyperref-generic} 
module (using \verb!\hypersetup!) in \texttt{hgb.sty}).



\section{Custom front pages using \emph{themes}}
\label{sec:Themes}

The content and structure of the front pages generated for the various thesis types
(\texttt{bachelor}, \texttt{master} \etc) may be easily customized to meet the specific
requirements of different institutions or departments.
Note that this kind of customization is typically done at the institutional level and not
by the individual author (student).


\subsection{Preconfigured themes}

The various front page arrangements are called \emph{themes} in the following.
Each theme is identified by a unique \texttt{\textit{themeID}} and associated with a 
particular style file (\latex\ package) named 
\texttt{hgbtheme-\textit{themeID}.sty}.
For example, the theme with \texttt{\textit{themeID}} = \texttt{default} is defined 
by the associated file
\begin{itemize}
  \item[] \texttt{hgbtheme-default.sty}.
\end{itemize}
%
Additional resources required by a theme (such as graphics files) must be named with the 
complete \emph{theme name} (\texttt{hgbtheme-\textit{themeID}})
as a prefix, for example,%
\footnote{This is to avoid file name conflicts since themes are part of the 
CTAN distribution and thus all theme-related files will eventually end up in a single, 
flat directory.}
\begin{itemize}
	\item[] \texttt{hgbtheme-fhooe26-logo.pdf}
\end{itemize}
%
which contains the logo graphics used by theme \texttt{fhooe26}.
The following themes are part of this package:
%
\begin{itemize}
	\item \texttt{default}: title page following the title page prototype for theses of the
		University of Applied Sciences Upper Austria (FH OÖ), with the logo at the top left
		and left-aligned text. This theme always uses the \emph{current} official brand
		(logo and name) of the institution (see Sec.~\ref{sec:ThemeBrand}).
	\item \texttt{fhooe26}: same layout as \texttt{default} with the FH OÖ brand pinned, \ie,
		it keeps the FH OÖ logo and name after future brand changes.
	\item \texttt{classic}: the former \texttt{default} theme (centered layout with the FH OÖ shield logo).
	\item \texttt{fhooe24}: title page with the FH OÖ cover background of 2024.
	\item \texttt{custom}: template for a custom theme (see Sec.~\ref{sec:CustomizingThemes}).
\end{itemize}

To \emph{use} a specific theme, option \texttt{theme=\textit{themeID}}
is added to the \verb!\documentclass! command, for example,
\begin{itemize}
	\item[] \verb!\documentclass[theme=default,...]{hgbthesis}!
\end{itemize}


\subsection{Brand: logo and institution name}
\label{sec:ThemeBrand}

Theme \texttt{default} keeps the \emph{brand}, \ie, the logo and the name of the institution,
separate from the layout of the title page. The brand is set by the command
\begin{itemize}
	\item[] \verb!\hgb@SetBrand[<logo height>]{<logo file>}{<name in German>}{<name in English>}!
\end{itemize}
%
where the optional logo height defaults to 28\,mm.
The institution's name is printed in the title page language, unless it is specified
explicitly by \verb!\institution{...}! in the document.
Brand themes like \texttt{fhooe26} load \texttt{default} as their parent theme and call
\verb!\hgb@SetBrand! afterwards, for example,
%
\begin{LaTeXCode}[numbers=none]
\hgb@UseTheme{default}	% parent theme (layout)
\hgb@SetBrand{my-logo.pdf}{Name der Hochschule}{Name of the institution}
\end{LaTeXCode}
%
The same approach can be used in a custom theme to adapt the title page to another institution.


\subsection{Customizing themes}
\label{sec:CustomizingThemes}

To customize the title page setup to their needs, authors (or administrators) should 
\emph{not} modify any of the standard theme files, since these may not be local but loaded
from a package distribution.
There are two recommended ways instead:

\subsubsection{Option 1: Adapt \texttt{hgbtheme-custom.sty}}

Copy the theme file \texttt{hgbtheme-custom.sty} (which is part of this distribution) to the 
\emph{main document directory} (if not there already) and open it in your LaTeX editor.
It defines \texttt{hgbtheme-default.sty} as the parent theme by the initial command
\begin{LaTeXCode}[numbers=none]
\hgb@UseTheme{default}
\end{LaTeXCode}
%
All commands defined in the parent theme are visible and may be redefined,
typically by \verb!\renewcommand! (see Sec.\ \ref{sec:ThemeStyleCommands}
for available commands).
To activate the associated theme simply use \texttt{custom} as the \texttt{\textit{themeID}}
in the associated document option, \ie,
%
\begin{LaTeXCode}[numbers=none]
\documentclass[theme=custom,...]{hgbthesis}
\end{LaTeXCode}
%
This is the simplest approach if only a \emph{single} custom theme is needed.

\subsubsection{Option 2: Create multiple custom themes}

Thesis administrators may find it useful to define \emph{multiple} custom themes for their
institution or department(s).
For this purpose, simply copy file \texttt{hgbtheme-default.sty} or \texttt{hgbtheme-custom.sty}
to a new file, \eg, \texttt{hgbtheme-physics.sty},%
\footnote{To be placed in the main document directory. Note the naming conventions!}
and modify it accordingly. The associated theme can then be invoked by
%
\begin{LaTeXCode}[numbers=none]
\documentclass[theme=physics,...]{hgbthesis}
\end{LaTeXCode}
%
An error will be raised if the associated \texttt{.sty} file cannot be found.



\subsection{Commands and variables available to custom theme styles}
\label{sec:ThemeStyleCommands}

Any theme style must at least define the command
\begin{itemize}
  \item[] \verb!\hgb@MakeFrontPages{..}!
\end{itemize}
which is pre-defined (as throwing an error) in \texttt{hgbthesis.cls}
and invoked by \verb!\maketitle!.
If \texttt{default} is used as the parent theme, \verb!\hgb@MakeFrontPages! 
is already set up properly and redefinition is optional, \ie,
%
\begin{LaTeXCode}[numbers=none]
\hgb@UseTheme{default}		% parent theme
\renewcommand{\hgb@MakeFrontPages}{..}	% optional
\end{LaTeXCode}
%

\noindent
Moreover, the following macros and variables are assured to be available for defining custom themes
(see \texttt{hgbtheme-default.sty} for their typical usage).
These are defined in \texttt{hgbthesis.cls} and should \emph{not} be redefined:
%
\begin{itemize}
  \item[] \verb!\hgb@Author!
  \item[] \verb!\hgb@Title!
  \item[] \verb!\hgb@SubTitle!
  \item[] \verb!\hgb@Institution!
  \item[] \verb!\hgb@ProgramType!
  \item[] \verb!\hgb@ProgramName!
  \item[] \verb!\hgb@AcademicDegree! (empty if not specified)
  \item[] \verb!\hgb@PlaceOfStudy!
\medskip
	\item[] \verb!\hgb@ThesisName!
	\item[] \verb!\hgb@ProposalName!
\medskip
  \item[] \verb!\hgb@AdvisorCount!
  \item[] \verb!\hgb@getAdvisorRole{<number>}!
  \item[] \verb!\hgb@getAdvisorName{<number>}!
\medskip
  \item[] \verb!\hgb@MainLanguage!
  \item[] \verb!\hgb@TitleLanguage!
  \item[] \verb!\hgb@TitlePageFont!
  %\item[] \verb!\hgb@LogoFile!
\medskip 
  %\item[] \verb!\hgb@Advisor!
  \item[] \verb!\hgb@SubmissionYear!
  \item[] \verb!\hgb@SubmissionMonth!
  \item[] \verb!\hgb@SubmissionDay!
  \item[] \verb!\hgb@GetMonthName{<language>}{<monthnumber>}!
\medskip
	\item[] \verb!\hgb@License!
\medskip
  \item[] \verb!hgb@IsProposal! (\emph{boolean}, without \verb!\!)
\end{itemize}

\noindent
Class \texttt{hgbthesis.cls} also defines a special \emph{hook}
\begin{itemize}
\item[] \verb!hgb@InitThemeHook!
\end{itemize}
for adding custom initialization code for the current theme in the form
%
\begin{LaTeXCode}[numbers=none]
\AddToHook{hgb@InitThemeHook}{<initialization code>},
\end{LaTeXCode}
%
typically placed at the beginning of the theme file. The collected code for
this hook is executed immediately \emph{before} \verb!\maketitle!.


In addition, theme \texttt{hgbtheme-classic.sty} defines the following
commands for generating individual title pages, which are inherited by theme \texttt{default}
(and thus by \texttt{fhooe26} and \texttt{custom}):
%
\begin{itemize}
  \item[] \verb!\hgb@MakeTitlePage! (redefined by \texttt{default})
  \item[] \verb!\hgb@MakeCompanyPage! (internship reports only)
  \item[] \verb!\hgb@MakeCopyrightPage!
  \item[] \verb!\hgb@MakeDeclarationPage!
\end{itemize}
%
Each of these may be redefined by inheriting themes.
Theme \texttt{default} additionally provides \verb!\hgb@SetBrand! (see Sec.~\ref{sec:ThemeBrand}).

\section{Package dependencies}

\begin{sloppypar}
The \texttt{hagenberg-thesis} package builds on the following \latex\ packages:\newline
\texttt{abstract}, 
\texttt{algorithm}, 
\texttt{algorithmicx}, 
\texttt{algpseudocodex}, 
\texttt{amsbsy}, 
\texttt{amsfonts}, 
\texttt{amsmath}, 
\texttt{amssymb}, 
\texttt{babel}, 
\texttt{biblatex}, 
\texttt{breakurl}, 
\texttt{caption}, 
\texttt{cmap}, 
\texttt{csquotes}, 
\texttt{datetime2}, 
\texttt{enumitem}, 
\texttt{epstopdf}, 
\texttt{exscale},
\texttt{extramarks},
\texttt{fancyhdr}, 
\texttt{float}, 
\texttt{fontenc},
\texttt{forloop},
\texttt{geometry}, 
\texttt{graphicx}, 
\texttt{hypcap}, 
\texttt{hyperref}, 
\texttt{ifpdf}, 
\texttt{inputenc}, 
\texttt{lengthconvert},
\texttt{lineno},
\texttt{listingsutf8}, 
\texttt{lmodern},
\texttt{marvosym}, 
\texttt{moreverb}, 
\texttt{overpic}, 
\texttt{pdfmanagement-testphase}, 
\texttt{pdfpages}, 
\texttt{pict2e}, 
\texttt{subdepth}, 
\texttt{titlesec}, 
\texttt{titling},
\texttt{tocbasic},
\texttt{url}, 
\texttt{upquote}, 
\texttt{verbatim}, 
\texttt{xcolor}, 
\texttt{xifthen},
\texttt{xstring},
\texttt{xspace}.
\end{sloppypar}

\end{document}
